% =====================================================================
% Paper 52 (standalone): Discrete spectral-triple realization of
% CFT_3-on-S^3 -- a non-holographic alternative to dS/CFT for boundary
% CFT data.
%
% Headline structural identification:
%   GeoVac sits in a third structural category in the
%   correspondence-physics landscape, distinct from Category I
%   (holographic: AdS/CFT, dS/CFT) and Category II (topological:
%   Chern-Simons/WZW). Category III is the spectral-triple
%   discretization category, with the framework's S^3 substrate
%   reproducing boundary CFT_3 data via propinquity convergence
%   (Paper 38), without any bulk gravitational dual.
%
% Substantive new structural content:
%   1. Three-category taxonomy of correspondence-physics relations
%   2. Articulation of GeoVac's actual position in the landscape
%   3. Connection of propinquity convergence (Paper 38) to the
%      "correspondence" mechanism role
%   4. Explicit distinction from both AdS/CFT and dS/CFT
%
% Fourteenth math.OA standalone in the GeoVac series.
% Siblings: Papers 38, 39, 40, 42, 43, 44, 45, 46, 47, 48, 49, 50.
%
% Status: first draft, May 2026, post-Sprint dS/CFT scoping
% (2026-05-29). Built from:
%   debug/sprint_ds_cft_scoping_memo.md
%   debug/geovac_correspondence_position_memo.md
%   debug/thread_3_synthesis_memo.md
%   papers/group1_operator_algebras/paper_50_cft3_partition_function.tex
%     (sec:catalogue, particularly the structural-distinction paragraph)
%
% Honest scope: framework-positioning paper, NOT a new holographic
% correspondence. The paper does NOT claim GeoVac is a Theory of
% Everything, a refutation of dS/CFT, or a new bulk-boundary duality.
% =====================================================================
\documentclass[11pt,a4paper]{article}

\usepackage[T1]{fontenc}
\usepackage[utf8]{inputenc}
\usepackage{amsmath, amssymb, amsthm, mathrsfs, mathtools}
\usepackage{geometry}
\geometry{margin=1in}
\usepackage{hyperref}
\hypersetup{colorlinks=true, linkcolor=blue!50!black, citecolor=blue!50!black, urlcolor=blue!50!black}
\usepackage{graphicx}
\usepackage{booktabs}
\usepackage[numbers,sort&compress]{natbib}
\usepackage{xcolor}

\theoremstyle{plain}
\newtheorem{theorem}{Theorem}[section]
\newtheorem{lemma}[theorem]{Lemma}
\newtheorem{proposition}[theorem]{Proposition}
\newtheorem{corollary}[theorem]{Corollary}
\theoremstyle{definition}
\newtheorem{definition}[theorem]{Definition}
\newtheorem{remark}[theorem]{Remark}
\newtheorem{example}[theorem]{Example}

\newcommand{\nmax}{n_{\max}}
\newcommand{\sthree}{S^{3}}
\newcommand{\Hilb}{\mathcal{H}}
\newcommand{\Acal}{\mathcal{A}}
\newcommand{\Tcal}{\mathcal{T}}
\newcommand{\R}{\mathbb{R}}
\newcommand{\C}{\mathbb{C}}
\newcommand{\N}{\mathbb{N}}
\newcommand{\Z}{\mathbb{Z}}
\newcommand{\Q}{\mathbb{Q}}
\newcommand{\Tr}{\mathrm{Tr}}
\newcommand{\DGV}{D_{\mathrm{GV}}}
\newcommand{\DCH}{D_{\mathrm{CH}}}
\newcommand{\Kalpha}{K_{\alpha}^{W}}
\newcommand{\rhoW}{\rho_{W}}
\newcommand{\KPS}{\text{KPS}}

\title{Discrete spectral-triple realization of\\
CFT$_{3}$-on-$\sthree$ partition function data:\\
a non-holographic alternative to dS/CFT}
\author{J.~Loutey\thanks{Independent researcher.
\texttt{jloutey@gmail.com}. This work was developed in an
AI-augmented agentic workflow with Anthropic's Claude.
Implementation, internal memos, and bibliographic collation were
performed collaboratively; all theorems, design choices, and editorial
decisions are the author's responsibility.}}
\date{May 2026 (preprint)}

\begin{document}
\maketitle

\begin{abstract}
The boundary CFT$_{3}$-on-$\sthree$ partition function data of the
GeoVac discrete spectral triple --- including bit-exact matches to the
Klebanov--Pufu--Safdi free energies of the conformally coupled scalar
and the free Weyl Camporesi--Higuchi Dirac~\cite{paper50,
klebanov_pufu_safdi2011} --- is reproduced without any bulk
gravitational theory. The framework's $\sthree$ substrate occupies
the geometric position that Euclidean dS/CFT~\cite{maldacena2002,
strominger2001} assigns to its boundary CFT$_{3}$, but the
construction lacks the bulk dS$_{4}$ (or its Euclidean partner $S^{4}$),
the Hartle--Hawking wavefunctional, and the holographic dictionary.
We argue that GeoVac sits in a structurally distinct \emph{third
category} of correspondence-physics relations, separated from both
\emph{Category I} (holographic: AdS/CFT, dS/CFT) and \emph{Category II}
(topological: Chern--Simons/WZW). \emph{Category III} is the
spectral-triple-discretization category; the framework's analog of the
holographic correspondence is the van~Suijlekom state-space
Gromov--Hausdorff convergence theorem~\cite{paper38} operating between two
operator-algebraic discretization levels (discrete spectral triple at
finite cutoff $\nmax$ and continuum spectral triple on round
$\sthree$) at convergence rate $C_{3} \gamma_{\nmax}$. This is not a
new holographic correspondence and not a refutation of dS/CFT; it is
a structurally distinct realization that produces the same boundary
CFT data via a different mechanism. We articulate the taxonomy, place
GeoVac within it, and identify the operational consequences of the
Category III reading for the framework's existing bit-exact boundary
CFT$_{3}$ results and for the modular structure of the wedge KMS
state~\cite{paper42, paper43}.
\end{abstract}

\tableofcontents

% =====================================================================
\section{Introduction}
\label{sec:intro}
% =====================================================================

\subsection{The dS/CFT question}

Strominger's dS/CFT proposal~\cite{strominger2001} posited a duality
between gravity on $(d+1)$-dimensional de Sitter spacetime and a
Euclidean CFT$_{d}$ on its conformal boundary. For $d = 3$, the
proposal places a CFT$_{3}$ on the future timelike infinity $I^{+}$
of dS$_{4}$, which is geometrically a copy of $\sthree$. Maldacena's
Euclidean sharpening~\cite{maldacena2002} via Wick rotation makes
the boundary identification precise: Euclidean dS$_{4}$ is the round
$S^{4}$, the equatorial $\sthree$ hosts the boundary CFT$_{3}$, and
the Hartle--Hawking wavefunctional of dS$_{4}$ at $I^{+}$ equals the
boundary CFT$_{3}$ partition function:
\begin{equation}
\Psi_{\mathrm{dS}_{4}}[\phi_{0}] \;=\; Z_{\mathrm{CFT}_{3}}[\phi_{0}].
\label{eq:dS_CFT_partition}
\end{equation}

The conjectural Eq.~\eqref{eq:dS_CFT_partition} relates a 4D bulk
quantity (a path integral over geometries with appropriate boundary
conditions) to a 3D boundary quantity (a CFT partition function on
$\sthree$).

\subsection{The GeoVac substrate as boundary geometry}

The GeoVac framework~\cite{paper7, paper38} constructs a discrete
spectral triple
$\Tcal_{\nmax} = (\Acal_{\nmax}, \Hilb_{\nmax}, \DGV)$ on the
Fock-projected $\sthree$. This $\sthree$ is the same manifold ---
same topology, same SO(4) symmetry, same Camporesi--Higuchi
spinor bundle~\cite{camporesi_higuchi1996} --- that
Eq.~\eqref{eq:dS_CFT_partition} assigns as the host of the boundary
CFT$_{3}$. The framework's discrete spectral-zeta machinery on
$\Tcal_{\nmax}$ produces bit-exact reproductions of the standard
Klebanov--Pufu--Safdi (KPS)~\cite{klebanov_pufu_safdi2011} closed-form
F-theorem coefficients for both the conformally coupled scalar and
the free Weyl Dirac~\cite{paper50}. The KPS values are precisely
the boundary CFT$_{3}$ partition function data that
Eq.~\eqref{eq:dS_CFT_partition} would predict via Hartle--Hawking.

\subsection{The structural question}

The GeoVac framework produces the boundary CFT$_{3}$ data
\emph{without invoking dS$_{4}$}: there is no bulk de Sitter
geometry in the framework's existing infrastructure, no Hartle--Hawking
wavefunctional construction, no functional integral over 4D
geometries, and no holographic dictionary mapping bulk fields to
boundary operators. The boundary CFT$_{3}$ values emerge directly
from spectral asymptotics of the discrete spectral triple ---
specifically, from the spectral-zeta derivative $-\tfrac{1}{2}
\zeta'_{\Delta}(0)$ evaluated on the round-$\sthree$ continuum limit
of $\DGV$~\cite{paper50}.

This raises a structural question: \emph{what is the framework's
position in the correspondence-physics landscape?} Is GeoVac a
non-standard realization of dS/CFT? An alternative to dS/CFT? A
different correspondence entirely?

\subsection{This paper's contribution}

We argue that GeoVac sits in a third structural category, distinct
from both AdS/CFT (Category I) and dS/CFT (Category I) on one hand,
and from topological constructions like Chern--Simons/WZW
(Category II) on the other. Category III is the
spectral-triple-discretization category, with two distinguishing
features:
\begin{enumerate}
\item There is no bulk gravitational theory; the framework's primary
object is a discrete spectral triple at finite cutoff, not a 4D
geometry to be holographic from.
\item The ``correspondence'' is between two operator-algebraic
objects (the discrete spectral triple at finite cutoff and the
continuum spectral triple on round $\sthree$), realized via
van~Suijlekom state-space Gromov--Hausdorff convergence~\cite{paper38} at rate
$C_{3} \gamma_{\nmax} \to 0$.
\end{enumerate}

The Connes--van Suijlekom truncation of the Toeplitz algebra on the
unit circle $S^{1}$~\cite{connes_vs2021} is the structural precedent
for Category III; GeoVac's $\sthree$ construction is the natural
non-abelian extension via Paper 38's five-lemma state-space
Gromov--Hausdorff convergence proof.

We do not claim that GeoVac is a new holographic correspondence
(it cannot be: there is no bulk). We do not claim it refutes
dS/CFT (the boundary data is consistent with what dS/CFT would
predict). We argue instead that GeoVac provides a structurally
distinct \emph{realization} of the boundary CFT$_{3}$-on-$\sthree$
partition function data, via a discrete spectral-triple mechanism that
operates differently from holographic reconstruction.

\subsection{Outline}

Section~\ref{sec:taxonomy} articulates the three-category taxonomy.
Section~\ref{sec:geovac_substrate} reviews the GeoVac substrate at
the level needed for the structural argument.
Section~\ref{sec:boundary_data} summarizes the bit-exact KPS matches
from Paper 50~\cite{paper50}.
Section~\ref{sec:propinquity_mechanism} identifies state-space GH
convergence (Paper 38) as the framework's analog of the holographic
correspondence.
Section~\ref{sec:modular_structure} addresses the wedge KMS modular
structure and its place in the boundary side of correspondence
theories.
Section~\ref{sec:comparison} explicitly compares the framework to
AdS/CFT and dS/CFT.
Section~\ref{sec:open} lists open questions.

% =====================================================================
\section{The three-category taxonomy of correspondence-physics relations}
\label{sec:taxonomy}
% =====================================================================

We organize the landscape of ``$X$/CFT-like'' relations into three
structural categories.

\subsection{Category I --- Holographic (bulk gravity $\leftrightarrow$ boundary CFT)}

\textbf{Examples.} AdS/CFT~\cite{maldacena1997}, dS/CFT~\cite{strominger2001,
maldacena2002}, Kerr/CFT, the flat-space/BMS analog.

\textbf{Structural ingredients.}
\begin{itemize}
\item A bulk gravitational theory in $d + 1$ dimensions (anti-de
Sitter, de Sitter, Kerr, flat space, etc.).
\item A boundary theory in $d$ dimensions, typically a CFT.
\item A holographic dictionary mapping bulk fields $\phi_{\rm bulk}$ to
boundary operators $\mathcal{O}_{\rm boundary}$ via an asymptotic
relation $\phi_{\rm bulk}(x, z) \sim z^{\Delta}\mathcal{O}_{\rm boundary}(x)$
as $z \to 0$.
\item A partition-function or wavefunction identification:
$Z_{\rm bulk}[\phi_{0}] = Z_{\rm boundary}[\phi_{0}]$ (AdS/CFT) or
$\Psi_{\rm bulk}[\phi_{0}] = Z_{\rm boundary}[\phi_{0}]$ (dS/CFT).
\end{itemize}

\textbf{Operationally.} Boundary CFT data is RECONSTRUCTED from bulk
gravity by holographic computation. Equivalently, bulk gravity is
computed from boundary CFT via reconstruction techniques (RT/HRT
minimum surfaces, JLMS bulk-modular reconstruction).

\subsection{Category II --- Topological (discrete bulk $\leftrightarrow$ boundary state)}

\textbf{Examples.} Chern--Simons / WZW boundary~\cite{witten1989},
state-sum models, tensor networks and MERA.

\textbf{Structural ingredients.}
\begin{itemize}
\item A bulk in $d + 1$ dimensions, typically with a discrete or
topological structure (lattice, simplicial complex, tensor network).
\item A boundary theory on a $d$-manifold, derived from bulk
topological data.
\item An identification of the boundary partition function with a
bulk topological invariant.
\end{itemize}

\textbf{Operationally.} Boundary observables emerge as topological
invariants of the bulk discrete structure. The bulk is finite or
discrete; the boundary CFT (or its analog) emerges in a continuum
limit of the discrete bulk data.

\subsection{Category III --- Spectral-triple discretization (substrate $\leftrightarrow$ continuum CFT)}

\textbf{Examples.} Connes--van Suijlekom truncation of the Toeplitz
algebra on $S^{1}$~\cite{connes_vs2021}; the GeoVac $\sthree$
extension via Paper~38's state-space GH convergence proof~\cite{paper38}; the present
paper's KPS-reproducing construction~\cite{paper50}.

\textbf{Structural ingredients.}
\begin{itemize}
\item A discrete spectral triple
$\Tcal_{\nmax} = (\Acal_{\nmax}, \Hilb_{\nmax}, D_{\nmax})$ at
finite cutoff.
\item A continuum spectral triple
$\Tcal_{\infty} = (\Acal_{\infty}, \Hilb_{\infty}, D_{\infty})$ in
the limit.
\item A propinquity-style metric $\Lambda$ on the space of spectral
triples.
\item Convergence: $\Lambda(\Tcal_{\nmax}, \Tcal_{\infty}) \to 0$ at
proven rate (Paper 38: $C_{3} \cdot \gamma_{\nmax}$ with
$\gamma_{\nmax} \sim (4/\pi) \log\nmax / \nmax$).
\item Continuum CFT$_{d}$ partition function data is the SPECTRAL
content of the limit triple, extracted via zeta-function derivatives
(KPS-style).
\end{itemize}

\textbf{Operationally.} Boundary CFT data is COMPUTED DIRECTLY from
spectral asymptotics of the discrete substrate, not reconstructed from
any bulk. There is no ``bulk'' --- the discrete substrate is the
framework's intrinsic object, and the continuum CFT is what the
substrate looks like at the state-space GH limit.

\subsection{The three categories at a glance}

\begin{center}
\begin{tabular}{lll}
\toprule
Category & Bulk object & Correspondence mechanism \\
\midrule
I (Holographic) & Bulk gravity in $d+1$ dimensions & Holographic dictionary; reconstruction \\
II (Topological) & Bulk discrete / topological in $d+1$ dimensions & State-sum; topological invariant \\
III (Spectral) & Discrete spectral triple at finite cutoff & State-space GH convergence; spectral asymptotics \\
\bottomrule
\end{tabular}
\end{center}

The three categories are not mutually exclusive in principle: one
could imagine constructions that combine, e.g., a holographic bulk
with a topological boundary lattice. We do not survey such hybrid
constructions here; we focus on the structural placement of GeoVac
as a Category III paradigm.

% =====================================================================
\section{The GeoVac substrate}
\label{sec:geovac_substrate}
% =====================================================================

The framework's primary geometric object is the discrete spectral
triple
\begin{equation}
\Tcal_{\nmax} = (\Acal_{\nmax}, \Hilb_{\nmax}, \DGV)
\label{eq:geovac_triple}
\end{equation}
on the Fock-projected~\cite{paper7} round $\sthree$, parameterized by
a cutoff $\nmax \in \N$. The Hilbert space
$\Hilb_{\nmax} = \bigoplus_{n=1}^{\nmax} \Hilb_{n}$ is the truncation
of the full $L^{2}$ spinor bundle on $\sthree$ at the $\nmax$-th
shell; the algebra $\Acal_{\nmax}$ is the (algebraic) Fock multiplier
algebra at that cutoff; $\DGV$ is the discrete Camporesi--Higuchi
Dirac operator with eigenvalues $|\lambda_{n}| = n + 3/2$ and
multiplicities $g_{n}^{\rm Dirac} = 2(n+1)(n+2)$ on the spinor bundle.

The continuum limit $\Tcal_{\infty}$ is the round-$\sthree$
Camporesi--Higuchi spectral triple on the full spinor $L^{2}$, with
the same eigenvalue structure extended to all $n \in \N$.

The substrate has no ``higher-dimensional'' geometry. There is no
4D bulk, no internal $z$-coordinate, no holographic radial direction.
The $\sthree$ is the framework's natural geometric arena, intrinsic to
the Fock projection of the Coulomb spectrum~\cite{paper7}.

\begin{remark}[Why no bulk]
\label{rem:no_bulk}
A natural question is whether the GeoVac framework could be
\emph{extended} with a 4D bulk geometry (Lorentzian dS$_{4}$ or
Euclidean $S^{4}$). The substrate's existing infrastructure does not
contain any such geometry: Paper 50~\cite{paper50} \S 6 documents that
the framework's codebase has no AdS$_{4}$ / H$^{4}$ / dS$_{4}$ / $S^{4}$
machinery, and Sprint RH-B\footnote{Sprint/track designations are anchors into the project's internal chronicle (\texttt{CHANGELOG.md} in the repository).} (2026-04-17)~\cite{paper29} closed the
$\sthree \to H^{3}$ Wick continuation as a clean structural dead end
(the continuation supplies no natural discrete subgroup
$\Gamma \subset \mathrm{SO}(3, 1)$ selectable from framework
invariants). Importing a 4D bulk would be a major structural
extension, beyond sprint scale, outside the present paper's scope and would, in any case,
be an \emph{extension} of the framework rather than something
already present.
\end{remark}

% =====================================================================
\section{The boundary CFT$_{3}$ partition function data}
\label{sec:boundary_data}
% =====================================================================

Paper 50~\cite{paper50} establishes that the GeoVac substrate
reproduces standard boundary CFT$_{3}$-on-$\sthree$ partition function
data bit-exactly. We summarize the headline results here for
self-containment.

\subsection{Conformally coupled scalar}

\begin{theorem}[Paper 50 \cite{paper50}, Theorem~3.1]
\label{thm:scalar_match}
On the truncated GeoVac spectral triple $\Tcal_{\nmax}$, the
spectral-zeta derivative of the conformally coupled scalar Laplacian
on round unit $\sthree$ satisfies
\begin{equation}
F_{s}^{\sthree} \;:=\; -\tfrac{1}{2}\, \zeta'_{\Delta_{\rm conf}}(0)
\;=\; \frac{\log 2}{8} - \frac{3 \zeta(3)}{16 \pi^{2}}.
\label{eq:Fs_S3}
\end{equation}
This matches the Klebanov--Pufu--Safdi closed
form~\cite{klebanov_pufu_safdi2011} at $61+$ digits and admits the
PSLQ integer relation $[8, 2, -3]$ across the basis
$\{\log 2, \zeta(3)/\pi^{2}\}$.
\end{theorem}

\subsection{Free Weyl Dirac}

\begin{theorem}[Paper 50 \cite{paper50}, Theorem~3.2]
\label{thm:dirac_match}
On the truncated GeoVac spectral triple $\Tcal_{\nmax}$, the
spectral-zeta derivative of the free Weyl Camporesi--Higuchi Dirac
on round unit $\sthree$ satisfies
\begin{equation}
F_{D}^{\sthree} \;:=\; -\tfrac{1}{2}\, \zeta'_{\DCH}(0)
\;=\; \frac{\log 2}{4} + \frac{3 \zeta(3)}{8 \pi^{2}}.
\label{eq:FD_S3}
\end{equation}
This matches the KPS closed form symbolically (sympy-exact).
\end{theorem}

\subsection{Orthogonal master Mellin engine decomposition}

The scalar and Dirac F-values combine linearly to extract orthogonal
master Mellin engine sectors~\cite{paper50, paper18}:
\begin{equation}
F_{D} + 2 F_{s} \;=\; \frac{\log 2}{2}
\qquad (\text{M2 only, all $\zeta(3)$ cancels})
\label{eq:M2_isolation}
\end{equation}
\begin{equation}
F_{D} - 2 F_{s} \;=\; \frac{3 \zeta(3)}{4 \pi^{2}}
\qquad (\text{M3 only, all $\log 2$ cancels}).
\label{eq:M3_isolation}
\end{equation}

This is the framework's first verification of the master Mellin engine
on non-GeoVac-internal physics observables~\cite{paper50, paper32}.

\subsection{The data the framework reproduces}

Eqs.~\eqref{eq:Fs_S3}--\eqref{eq:M3_isolation} are standard CFT$_{3}$
data. Eq.~\eqref{eq:Fs_S3} is the F-theorem
coefficient~\cite{klebanov_pufu_safdi2011} that monotonically decreases
along renormalization-group flows of 3D CFTs. The Maldacena
identification~\eqref{eq:dS_CFT_partition} would predict exactly these
values via the holographic computation of the Hartle--Hawking
wavefunctional of dS$_{4}$ at $I^{+}$. The framework reproduces them
without that computation.

% =====================================================================
\section{State-space GH convergence as the ``correspondence'' mechanism}
\label{sec:propinquity_mechanism}
% =====================================================================

The framework's analog of the holographic correspondence is
van~Suijlekom state-space Gromov--Hausdorff convergence between the
discrete spectral triple at finite cutoff and its continuum limit.

\subsection{Statement of the convergence theorem}

\begin{theorem}[Paper 38 \cite{paper38}; WH1 PROVEN]
\label{thm:propinquity}
The discrete GeoVac spectral triples $\Tcal_{\nmax}$ on the
Fock-projected $\sthree$ converge to the continuum round-$\sthree$
Camporesi--Higuchi spectral triple $\Tcal_{\infty}$ in the
van~Suijlekom state-space Gromov--Hausdorff sense:
\begin{equation}
\Lambda(\Tcal_{\nmax}, \Tcal_{\infty})
\;\le\; C_{3} \cdot \gamma_{\nmax}
\;\to\; 0 \quad \text{as } \nmax \to \infty,
\label{eq:propinquity_rate}
\end{equation}
where $C_{3} = 1$ asymptotically (Lipschitz comparison constant) and
$\gamma_{\nmax} = O(\log \nmax / \nmax)$ is the central-Fej\'er
moment with universal asymptotic rate constant $4/\pi$.
\end{theorem}

The proof is the five-lemma chain $\mathrm{L1}'$-$\mathrm{L2}$-$\mathrm{L3}$-$\mathrm{L4}$-$\mathrm{L5}$
of Paper 38~\cite{paper38}, generalized to all simple compact Lie
groups in Paper 40~\cite{paper40} (for general $G$ conditional on a
per-group spin-window decomposition, verified for $\mathrm{SU}(2)$).

\subsection{Why this is the ``correspondence'' mechanism}

The state-space GH convergence relates two operator-algebraic objects
(the discrete and continuum spectral triples) by a quantitative
metric, with proven convergence rate. The boundary CFT$_{3}$
partition function data of Section~\ref{sec:boundary_data} is
extracted from the continuum spectral triple $\Tcal_{\infty}$ via the
zeta-function derivative $-\tfrac{1}{2}\zeta'(0)$; this derivative is
PRESERVED by the state-space GH limit because
\begin{itemize}
\item the spectral data of $D_{\nmax}$ converges to that of
$D_{\infty}$ in the operator-norm sense (per Berezin reconstruction,
Paper 38 L4);
% C7: state-space GH (van Suijlekom), not Latrémolière propinquity
\item the spectral-zeta function converges term-by-term and uniformly
on the strip $\mathrm{Re}(s) > d/2$;
\item the meromorphic continuation to $s = 0$ commutes with the
state-space GH limit by standard analytic-continuation arguments.
\end{itemize}

In holographic terminology, this is the framework's ``dictionary'':
the discrete spectral triple at finite cutoff is the framework's
``bulk-like'' object (more precisely, its substrate), and the
boundary CFT$_{3}$ data is what that object looks like at the
state-space GH limit. There is no separate ``bulk theory'' to translate
from; the spectral triple at finite cutoff is itself the framework's
primary object.

\subsection{Difference from holographic reconstruction}

Holographic reconstruction recovers boundary CFT data from bulk
gravity by integrating field-theoretic correlators against bulk
modes. The boundary CFT exists in the same dimensional setting as the
bulk's asymptotic boundary, and the dictionary translates between two
distinct physical theories.

State-space GH convergence recovers continuum CFT data from a discrete
substrate by taking the continuum limit of the substrate's spectral
asymptotics. Both objects (discrete and continuum spectral triples)
live in the same dimensional setting (both are spectral triples on
$\sthree$); the dictionary translates between two operator-algebraic
discretization levels, not between two physical theories of different
dimension.

This is a structurally distinct mechanism. It is appropriate to call
both a ``correspondence'' in the abstract sense (each relates
mathematical structures in a way that yields physical data), but the
specific nature of the correspondence differs.

% =====================================================================
\section{Wedge KMS modular structure}
\label{sec:modular_structure}
% =====================================================================

A second structural feature of the GeoVac framework that overlaps
with holographic theories is the wedge KMS modular
structure~\cite{paper42, paper43, paper50}.

\subsection{Wedge KMS state on the framework}

Paper 50~\cite{paper50} defines a wedge KMS state $\rhoW = e^{-\Kalpha} / Z$
on a hemispheric wedge of the $\sthree$ substrate, where $\Kalpha$ is
the modular Hamiltonian associated with the boost generator of the
wedge. The state $\rhoW$ has the structural properties of the
Bisognano--Wichmann vacuum~\cite{bisognano_wichmann1976}.

Paper 42~\cite{paper42} establishes the four-witness Wick-rotation
theorem: the modular structure of $\rhoW$ corresponds bit-exactly with
the thermal-state readings of Hawking, Sewell, Bisognano--Wichmann,
and Unruh at finite cutoff $\nmax$. The Tomita--Takesaki modular flow
on the wedge KMS state is a literal identification at the
operator-system level (Riemannian).

\subsection{Connection to holographic theories}

Any holographic correspondence involving thermal states
(AdS-Schwarzschild $\leftrightarrow$ thermal CFT, dS-Schwarzschild
$\leftrightarrow$ static patch CFT, BTZ $\leftrightarrow$ thermal
CFT on $S^{1} \times \R$) has the SAME modular-flow structure on its
boundary CFT side. The four-witness theorem of Paper 42 explicitly
identifies Hawking, Sewell, Bisognano--Wichmann, and Unruh as four
readings of one modular-flow structure.

GeoVac therefore shares the boundary-modular-flow content with all
holographic-correspondence theories. The Connes--Rovelli thermal-time
identification of Paper 49~\cite{paper49} makes this concrete:
thermal time is the modular flow, regardless of whether one came at
it from a bulk-gravity calculation or a direct boundary-CFT analysis.

The structural distinction we are drawing is therefore:
\begin{itemize}
\item The boundary-modular-flow content is SHARED across Category I
and Category III.
\item The bulk content distinguishes them: Category I has a bulk
gravitational theory; Category III has a discrete spectral-triple
substrate.
\end{itemize}

% =====================================================================
\section{Comparison to AdS/CFT and dS/CFT}
\label{sec:comparison}
% =====================================================================

We summarize the comparison in a cross-table.

\begin{center}
\begin{tabular}{lll}
\toprule
Theory & Bulk object & Correspondence mechanism \\
\midrule
AdS$_{5}$/CFT$_{4}$~\cite{maldacena1997} & AdS$_{5}$ IIB SUGRA & Holographic dictionary \\
AdS$_{4}$/CFT$_{3}$ (ABJM) & AdS$_{4}$ M-theory & Holographic dictionary \\
dS$_{4}$/CFT$_{3}$~\cite{strominger2001, maldacena2002} & dS$_{4}$ & Hartle--Hawking wavefunctional \\
Kerr/CFT & Kerr & Asymptotic-symmetry holography \\
Chern--Simons/WZW~\cite{witten1989} & 3D topological & State-sum \\
Connes--vS $S^{1}$~\cite{connes_vs2021} & Discrete spectral triple on $S^{1}$ & Propinquity (Toeplitz $S^{1}$) \\
\textbf{GeoVac} (Papers~\cite{paper38, paper50}) & \textbf{Discrete spectral triple on $\sthree$} & \textbf{State-space GH convergence at rate $4/\pi$} \\
\bottomrule
\end{tabular}
\end{center}

\subsection{What GeoVac shares with dS/CFT}

\begin{itemize}
\item The same boundary host manifold: the Euclidean dS/CFT
identification~\eqref{eq:dS_CFT_partition} places the boundary
CFT$_{3}$ on $\sthree$; GeoVac's substrate is on the same $\sthree$.
\item The same boundary CFT$_{3}$ partition function data:
Eqs.~\eqref{eq:Fs_S3}, \eqref{eq:FD_S3} reproduce values that
Hartle--Hawking would predict.
\item The same wedge KMS / modular-flow structure on the boundary side
(four-witness theorem, Paper 42~\cite{paper42}).
\end{itemize}

\subsection{What GeoVac does not share with dS/CFT}

\begin{itemize}
\item GeoVac has no bulk dS$_{4}$ (Lorentzian) or $S^{4}$ (Euclidean).
\item GeoVac has no Hartle--Hawking wavefunctional construction.
\item GeoVac has no functional integral over 4D geometries.
\item GeoVac has no holographic dictionary mapping bulk fields to
boundary operators.
\end{itemize}

\subsection{The Category III alternative}

The framework's mechanism for producing the boundary CFT$_{3}$ data
is the state-space GH convergence of Theorem~\ref{thm:propinquity}: the
discrete spectral triple at finite cutoff approaches the continuum
round-$\sthree$ spectral triple in an explicit metric, with proven
convergence rate. The boundary CFT$_{3}$ data is the spectral-zeta
content of the continuum limit, and this content is preserved by the
state-space GH limit by standard analytic-continuation arguments. This
is structurally distinct from any holographic mechanism.

\begin{remark}[What this paper is NOT claiming]
\label{rem:not_claiming}
We explicitly do not claim:
\begin{enumerate}
\item GeoVac is a new holographic correspondence (no bulk $=$ no
holography in the standard sense).
\item GeoVac is a refutation of dS/CFT (the boundary data is
consistent; we just produce it differently).
\item GeoVac is a Theory of Everything (per the rhetoric rule of
the GeoVac project~\cite{paper7}, the framework does not assert
ontological priority of the spectral-triple description over
continuum quantum mechanics).
\item GeoVac proves Maldacena's identification~\eqref{eq:dS_CFT_partition}
(we have no access to the bulk side).
\end{enumerate}
What we do claim is that GeoVac produces the same boundary CFT$_{3}$
data via a structurally distinct mechanism that places the framework
in a third category, distinct from holographic correspondences.
\end{remark}

% =====================================================================
\section{Open questions}
\label{sec:open}
% =====================================================================

\paragraph{Q1.} The framework's bulk-side identifications corresponding
to Hartle--Hawking dS$_{4}$ (or its Euclidean $S^{4}$ partner), if any,
remain unreachable from the existing infrastructure. Importing
4D bulk machinery would be a major structural extension, beyond
sprint scale.
Whether such an extension is structurally compatible with the
framework's existing Coulomb / Fock-projection origin is open.

\paragraph{Q2.} The Connes--van Suijlekom Toeplitz $S^{1}$
truncation~\cite{connes_vs2021} is the abelian Category III precedent.
Paper 40~\cite{paper40} establishes the universal rate constant $4/\pi$
for all simple compact Lie groups; the state-space Gromov--Hausdorff
convergence itself is rigorous for $\mathrm{SU}(2)$ and, for general $G$,
conditional on a per-group spin-window decomposition (a named gap). Whether the
boundary CFT data for other groups (e.g., $\mathrm{SU}(3)$,
$\mathrm{Sp}(2)$, $G_{2}$) similarly reproduces standard CFT-on-group
partition functions is a natural extension.

\paragraph{Q3.} The boundary CFT$_{3}$ data of Eqs.~\eqref{eq:Fs_S3},
\eqref{eq:FD_S3} reproduces free-field KPS values. Whether the framework
extends to interacting CFT$_{3}$ partition function data is open. The
master Mellin engine~\cite{paper18, paper32} suggests that the
transcendental signature of any CFT$_{3}$ partition function decomposes
under M1/M2/M3, but mapping this decomposition to interacting CFT$_{3}$
data requires further analysis.

\paragraph{Q4.} A metric-level Lorentzian extension of Category III to
non-compact carriers does \emph{not} currently close:\ Papers
45~\cite{paper45} and 46 are descoped (the $K^{+}$ compression
annihilates the Lipschitz seminorm --- a degeneracy theorem, not a
Lorentzian convergence), and only Paper 47~\cite{paper47}'s
norm-resolvent arrow survives (partial). Whether any future Lorentzian
extension produces dS-thermal-state observables (Hawking radiation on
the discrete substrate, thermal density of states, etc.) that match
dS/CFT predictions is open.

% =====================================================================
\section*{Acknowledgments}
% =====================================================================

This paper consolidates findings from the post-v3.20.0 continuation
thread (2026-05-29) of the GeoVac development, building on Paper 50
(boundary CFT$_{3}$ partition function data) and Paper 38 (state-space
GH convergence). The framework-positioning argument was developed in
internal memos
\texttt{debug/sprint\_ds\_cft\_scoping\_memo.md}~\cite{debug:ds_cft_scoping}
and
\texttt{debug/geovac\_correspondence\_position\_memo.md}~\cite{debug:correspondence_position}.
The work was developed in an AI-augmented agentic workflow with
Anthropic's Claude.

% =====================================================================
\begin{thebibliography}{99}
% =====================================================================

\bibitem{paper7} J.~Loutey,
\emph{Paper 7: The Dimensionless Vacuum: Recovering the Schr\"odinger
Equation from Scale-Invariant Graph Topology},
GeoVac internal preprint, 2026.

\bibitem{paper18} J.~Loutey,
\emph{Paper 18: Spectral-Geometric Exchange Constants: Projecting
Discrete Spectra onto Continuous Manifolds},
GeoVac internal preprint, 2026.

\bibitem{paper29} J.~Loutey,
\emph{Paper 29: Finite-Size Graph-RH for the GeoVac Hopf Graphs: Ihara Zeta, Bound-Crossing, and a Scope Boundary on Selberg-on-Hydrogen},
GeoVac internal preprint, 2026.

\bibitem{paper32} J.~Loutey,
\emph{Paper 32: The GeoVac Spectral Triple: A Synthesis Paper Locking
the Framework into the Marcolli--van~Suijlekom Gauge-Network Lineage},
GeoVac internal preprint, 2026.

\bibitem{paper38} J.~Loutey,
\emph{Paper 38: State-space Gromov--Hausdorff convergence of truncated
Camporesi--Higuchi spectral triples on $S^3$},
GeoVac internal preprint, 2026.

\bibitem{paper40} J.~Loutey,
\emph{Paper 40: State-space Gromov--Hausdorff convergence of spectral truncations of compact Lie groups with bi-invariant metric},
GeoVac internal preprint, 2026.

\bibitem{paper42} J.~Loutey,
\emph{Paper 42: Tomita--Takesaki modular structure on truncated
$SU(2)$ spectral triples: four-witness Wick-rotation literal
identification at finite cutoff},
GeoVac internal preprint, 2026.

\bibitem{paper43} J.~Loutey,
\emph{Paper 43: Lorentzian extension of the four-witness Wick-rotation
theorem on truncated $S^3 \times \mathbb{R}$ spectral triples at finite
cutoff},
GeoVac internal preprint, 2026.

\bibitem{paper45} J.~Loutey,
\emph{Paper 45: Lorentzian propinquity on truncated $\mathrm{SU}(2)\times\mathrm{U}(1)_{T}$ Krein spectral triples: a degeneracy theorem},
GeoVac internal preprint, 2026.

\bibitem{paper47} J.~Loutey,
\emph{Paper 47: Norm-Resolvent Convergence to the Non-Compact
Lorentzian Krein Spectral Triple and the Two-Rate Hybrid Limit},
GeoVac internal preprint, 2026.

\bibitem{paper49} J.~Loutey,
\emph{Paper 49: OSLPLS Category and the Strong-Form Krein--Mondino--S\"amann
Bridge},
GeoVac internal preprint, 2026.

\bibitem{paper50} J.~Loutey,
\emph{Paper 50: Boundary CFT$_3$-on-$S^3$ observables on the truncated
GeoVac spectral triple: partition function, wedge KMS entropy, and
Casini--Huerta--Myers modular Hamiltonian},
GeoVac internal preprint, 2026.

\bibitem{camporesi_higuchi1996} R.~Camporesi and A.~Higuchi,
\emph{On the eigenfunctions of the Dirac operator on spheres and real
hyperbolic spaces},
J. Geom. Phys. \textbf{20} (1996) 1--18.

\bibitem{connes_vs2021} A.~Connes and W.~D.~van Suijlekom,
\emph{Spectral truncations in noncommutative geometry and operator systems},
Comm. Math. Phys. \textbf{383} (2021) 2021--2067, arXiv:2004.14115.

\bibitem{strominger2001} A.~Strominger,
\emph{The dS/CFT correspondence},
JHEP \textbf{10} (2001) 034, arXiv:hep-th/0106113.

\bibitem{maldacena1997} J.~Maldacena,
\emph{The large $N$ limit of superconformal field theories and supergravity},
Adv. Theor. Math. Phys. \textbf{2} (1998) 231--252, arXiv:hep-th/9711200.

\bibitem{maldacena2002} J.~Maldacena,
\emph{Non-Gaussian features of primordial fluctuations in single field inflationary models},
JHEP \textbf{05} (2003) 013, arXiv:astro-ph/0210603.

\bibitem{klebanov_pufu_safdi2011} I.~R.~Klebanov, S.~S.~Pufu, and B.~R.~Safdi,
\emph{F-Theorem without Supersymmetry},
JHEP \textbf{10} (2011) 038, arXiv:1105.4598.

\bibitem{bisognano_wichmann1976} J.~J.~Bisognano and E.~H.~Wichmann,
\emph{On the duality condition for quantum fields},
J. Math. Phys. \textbf{17} (1976) 303--321.

\bibitem{witten1989} E.~Witten,
\emph{Quantum field theory and the Jones polynomial},
Comm. Math. Phys. \textbf{121} (1989) 351--399.

\bibitem{debug:ds_cft_scoping} GeoVac internal,
\emph{Sprint dS/CFT scoping memo},
\texttt{debug/sprint\_ds\_cft\_scoping\_memo.md}, 2026-05-29.

\bibitem{debug:correspondence_position} GeoVac internal,
\emph{GeoVac correspondence position memo},
\texttt{debug/geovac\_correspondence\_position\_memo.md}, 2026-05-29.

\end{thebibliography}

\end{document}
